
\section{The Lutheran Free Church}

"Folkebladet" for April 28 and May 12, 1897.

"The United Church" has, through its twenty-year history, demonstrated that it has power neither to bring peace nor freedom into the Church. And yet it was in order to attain these two things that it was organized.

The people of the Church are left sitting with great disappointment and with empty hearts.

But the cause is neither fallen nor abandoned. In reality it is a step forward that there are already many—and more are becoming so—who acknowledge that peace and freedom in the Church cannot be attained by church politics and lawsuits.

Spirit and life and love must be added, if peace and freedom are to come among us.

It is this consciousness that has called forth the work for the Free Church, which is beginning to show small signs of stirring here and there among the congregations, and which also makes itself known through the proposed rules that are published today in Folkebladet."*

The thought is that in the Free Church the congregation should come up and the "association" come down. The intention is to attain the strength of cooperation by Spirit, not by law; by freedom, not by compulsion.

By rules are here meant "Guiding Principles" and "Rules for the Work." They are found on page {FIXME 49 ff}.

Two things, then, are necessary in the Free Church: 1) that the freedom and authority of the individual congregation be fully acknowledged and without any curtailment, and 2) that the individual congregation truly be filled and permeated by God’s Spirit.

Everything in the Free Church depends upon these two things: that the congregations become independent and spiritual. And in truth they cannot become the one without the other.

It is therefore only a pitiful and sorrowful spiritlessness and ignorance of godly things when one says that there is no need among us of any work for the Free Church; we already have it, because state and church are separated from one another.

Indeed? Is the state the only power that can make the Church unfree and that has made it so? If that were the case, then the Catholic Church would be an excellent Free Church. Surely it cannot be the intention to lead us back again into its nature? Or was it not precisely the Catholic bondage that Luther designated by the name: "The Babylonian Captivity of the Church"? No, the truth is this: the Lutherans preferred the dominion of the state in the Church to the hard tyranny of the Catholic Church and desired the protection of the state against this.

By the mere fact that we are out of the power of the state, we are no Free Church. The so-called "associations," or synods, exercise more dominion than the state would dare in this age.* But even if we did not have these worldly, lawsuit-mad associations over us, we would still be only a caricature of a free congregation, so long as sheer worldliness and death are not only found, but even rule in the congregations, as they do in many places.

Only where the Spirit of the Lord is, is there freedom. Therefore the work for the Free Church is the same thing as the work for the enlivening of the congregation.

In this work of awakening the Lord calls upon all believers to be laborers. All believers, both pastors and laypeople, both men and women, are called out into the Lord’s service to rescue life and save souls.

The old-fashioned view from the very worst days of Rationalism, that it is only the pastors who make up the Church’s labor-force and the Lord’s servants, is not sufficient in the great conflict between faith and unbelief into which we have been cast in our time. The principle that is to be applied is the universal duty of all believers to serve.

Therefore the thought is first and foremost to awaken believers to see their calling and to walk and work accordingly, so that they may become zealous to awaken the sleeping within and outside the congregation.

Only when the congregations are freed from the yoke of worldliness are they truly free.

But if the congregations now become so independent, can they then do any work in common? Will it not become with the congregations in the Free Church as with the so-called "outside" congregations, which the church bodies regarded with so much suspicion and in part with contempt?

That will not happen, if the believers in the congregations step forward in the work. They do indeed have the unity of the Spirit, and they know how to preserve it in the bond of peace. The voluntary laborers will have no difficulty in agreeing to work for the cause they love. And when each works for the cause according to his ability and gift, then it will become altogether different from the "outside" congregations, which often "stood outside" because they wished to escape the work, or because the bonds of the associations were too tight for them.

This Free Church work is therefore the only practical and real work of union—or rather work of peace.* Give room to all congregations, let each of them have its freedom and independence unimpaired, do not bind them together in bundles with the bonds of denominational constitutions, and then peace is very near to coming.

If the work has so great a goal, it must of course be content to begin quite small. It is never many who begin a spiritual struggle for the good of the whole. But we have the faith that it depends neither on the many nor on the few, but on the Lord’s Spirit, whether such work shall succeed.

But whether anyone is in an "association" or outside an "association," they are equally welcome to take part in the Free Church work, if they are agreed in this purpose: the enlivening and liberation of the congregations.

The proposal (for rules for a Lutheran Free Church) cannot be considered too carefully, nor criticized too sharply. For a Lutheran Free Church, in which the congregations truly enjoy that freedom of life and of the Spirit which belongs to them according to the New Testament, is something new, at least here in America.* It has been so, and it is still so rather generally to this day, that there is almost nowhere any possibility of joining a Lutheran congregation without at the same time also having to become a member of an ecclesiastical party.

But when a cause is so new,* then it is altogether certain that it is of great importance to proceed cautiously. And the right caution consists in laying as much weight as possible upon the congregation, that it may receive its right spiritual shape among us according to God’s Word; for there lies the Free Church’s strength and center of gravity, and if the congregation is allowed to come forth, then there will be "Free Church" of itself.

For "Free Church" means, so far as I understand anything of it, nothing other than this: that the congregation is allowed to be and to work in freedom* for the cause of God’s Kingdom, not as other people and congregations determine* for them, but as they determine themselves.

But if the congregations are truly free, that is to say, living and driven by the Spirit, then it will be a self-evident thing that those who are driven by one and the same Spirit will also find and understand one another, and mutually support and encourage one another. And here again the New Testament will be able to give us clear and sufficient guidance with respect to the connection of the free congregations with one another.\footnote{* See Professor Oftedal’s lecture, page FIXME 57 ff.}live the same divine life and are a chain of siblings among God’s children—be assured of it—they love one another, they understand one another, they have care for one another, they rejoice to see one another, and what interests the one cannot be a matter of indifference to the other. Therefore freedom will not produce strife and disorder, but precisely peace and order. Compulsion and injustice bring about strife and confusion, as we may see daily. Suppose we tried what freedom and love would accomplish! Or does anyone doubt that there will be better understanding where two brothers labor freely at one another’s side than where the one would threaten the other into being a slave?"

The right form of cooperation among congregations is common interest and full voluntariness. To this belongs that there be the same lively and spiritual cooperation among the congregations as the New Testament shows us in the Apostolic Age. We must visit one another, we must come together for great and small meetings; we must write letters, we must circulate papers and books.

In this way the Free Church will attain that mutual brotherly understanding which will make it possible for the work of the Church to take place in voluntariness, and that those who are most interested in the advancement of the cause may also have the most influence upon its course. A good cause will find the support that it can obtain by virtue of the interest it awakens. And more than that it surely has no right to demand.

If we could get free and spiritual congregations, then their cooperation would not be lacking, only that it would be the common labor of brothers and sisters in mutual love, not the wearying yoke of enforced denominational duty.

Georg Sverdrup

FOOTNOTE


